% Generated draft: requires mathematical and source review.
\documentclass[11pt]{article}
\usepackage[margin=25mm]{geometry}
\usepackage{fontspec}
\setmainfont{Latin Modern Roman}
\usepackage{amsmath,amssymb,mathtools,unicode-math}
\setmathfont{Latin Modern Math}
\usepackage{enumitem,needspace,xurl,hyperref}
\setlength{\parindent}{0pt}
\setlength{\parskip}{5pt}
\setlength{\emergencystretch}{3em}
\newcommand{\N}{\mathbb{N}}
\newcommand{\Z}{\mathbb{Z}}
\newcommand{\Q}{\mathbb{Q}}
\newcommand{\R}{\mathbb{R}}
\newcommand{\C}{\mathbb{C}}
\newcommand{\F}{\mathbb{F}}
\newcommand{\A}{\mathbb{A}}
\newcommand{\PP}{\mathbb{P}}
\newcommand{\E}{\mathbb{E}}
\newcommand{\eps}{\varepsilon}
\newcommand{\dd}{\,\mathrm{d}}
\begin{document}
\section*{UnsolvedMath ID: 700871}
\textbf{Source problem code:} catalogue entry
\subsection{Definitions and mathematical statement}
An \(\ell\)-group is a group \(G\) equipped with a partial order \(\leq\) for which every pair has a join \(x\vee y\) and a meet \(x\wedge y\), and which is invariant under left and right multiplication: for all \(a,b,x,y\in G\),
\[
x\leq y\implies axb\leq ayb.
\]
An \(\ell\)-homomorphism is a group homomorphism that preserves both joins and meets.

For an integer \(\eta\) with \(1<\eta<\infty\), let \(\mathcal F_{\eta}\) be the free \(\ell\)-group on generators \(x_1,\ldots,x_{\eta}\). Thus, for every \(\ell\)-group \(H\) and every choice \(h_1,\ldots,h_{\eta}\in H\), there is a unique \(\ell\)-homomorphism \(f:\mathcal F_{\eta}\to H\) satisfying \(f(x_i)=h_i\) for all \(i\).

The question of S. McCleary asks whether, for every such finite rank \(\eta\), the group \(\mathcal F_{\eta}\) is Hopfian in the category of \(\ell\)-groups; explicitly, whether every surjective \(\ell\)-homomorphism
\[
\varphi:\mathcal F_{\eta}\longrightarrow\mathcal F_{\eta}
\]
is injective. Equivalently, must every such \(\varphi\) satisfy \(\ker\varphi=\{1\}\)?
\subsection{Short English statement}
For every finite rank greater than one, must every surjective map from the free lattice-ordered group to itself that preserves the group and lattice operations be one-to-one? The supplied registry labels this question open; that status is not independently verified here.
\subsection{Sources}
\begin{itemize}
\item Local UnsolvedMath registry, permanent ID \(700871\), exact source code \texttt{catalogue entry}, title \emph{Kourovka Notebook Problem 12.52}; supplied complete source record.
\item S. McCleary, proposer, 1992. The supplied attribution identifies Problem 12.52 in \emph{The Kourovka Notebook (Unsolved Problems in Group Theory)}, 12th issue (1992).
\item The supplied background also identifies the 21st issue edition (2026), arXiv:1401.0300v48. The supplied collection description names E. I. Khukhro and V. D. Mazurov as editors. These bibliographic details are supplied attribution and have not been independently verified.
\end{itemize}
\end{document}
